\vspace{-0.5\baselineskip}
\section{Problem Definition and Background}
\vspace{-0.6\baselineskip}

\subsection{Inference-Time Reward Alignment}
\vspace{-0.4\baselineskip}
\label{sec:problem_definition}
Given a pretrained flow model that maps the source distribution, a standard Gaussian distribution $p_1$, into the data distribution $p_0$, our objective is to generate high-reward samples $\B{x}_0 \in \mathbb{R}^d$ from the pretrained flow model without additional training--a task known as inference-time reward alignment. We denote the given reward function as $r: \mathbb{R}^{d} \rightarrow \mathbb{R}$, 
which measures text alignment or user preference for a generated sample. 
% which is designed to measure text alignment or user preferences for the generated output.
% This function can capture various criteria, such as text alignment or user preferences, depending on the specific task. 
Following previous works~\cite{Korbak:2022RLDM,Uehara:2024Finetuning,Uehara:2024Bridging}, our objective can be formulated as finding the following target distribution:
{\small
    \begin{align}
        \label{eq:reward_max_obj}
        p^{*}_0 &= \argmax_{q} \  \mathbb{E}_{\B{x}_0 \sim q} \underbrace{\left[ r(\B{x}_0) \right]}_{\text{Reward}} -\beta \underbrace{\mathcal{D}_{\text{KL}} \left[ q \| p_0 \right]}_{\text{KL Regularization}},
    \end{align}
}%
% which covers samples with high reward while not deviating too much from the original data distribution $p(\B{x}_0)$, with its strength controlled by the hyperparameter $\beta$. 
which maximizes the expected reward while the KL divergence term prevents $p^{*}_0(\B{x}_0)$ from deviating too far from $p_0(\B{x}_0)$, with its strength controlled by the hyperparameter $\beta$. 
As shown in previous work~\cite{Rafailov2023:DPO}, the target distribution $p^{*}_0$ can be computed as:
{\small
\begin{align}
    \label{eq:target_distribution}
    p^{*}_0(\B{x}_0) &= \frac{1}{Z} p_0(\B{x}_0) \exp \left( \frac{r(\B{x}_0)}{\beta} \right),
\end{align}
}%
where $Z$ is a normalization constant. We present details in Appendix~\ref{sec:appendix_optimal_marginal_target}. However, sampling from the target distribution is non-trivial. 

% A notable approach for sampling from the target distribution is \emph{particle sampling}, which iteratively draws new particles from high-reward samples, focusing computation on promising regions. 
A notable approach for sampling from the target distribution is \emph{particle sampling}, which maintains a set of candidate samples—referred to as particles—and iteratively propagates high-reward samples while discarding lower-reward ones. 
% iteratively refines them by shifting toward higher-reward regions.
When combined with the denoising process of diffusion models, particle sampling can improve the efficiency of limited computational resources in inference-time scaling. 
In the next section, we review particle sampling methods used in diffusion models and, \emph{for the first time,} we explore insights for adapting them to flow models. 
% This improves efficiency by reducing unnecessary evaluations, driving its widespread use in inference-time scaling.


\mycomment{
1. 우리의 목적은 pretrained된 flow model로부터 high reward를 가진 샘플을 추가 학습 없이 생성해내는 것이며, 이는 reward alignment task라고도 불린다. 
2. 이 때 우리는 pretrained 된 flow 모델의 샘플 분포를 p_\theta, 주어진 reward model은 r(\cdot)이라고 한다. 
% 
3. Formally, 우리의 이러한 방식은 다음의 목적 함수로 정리된다: 
4. 여기서 KL divergence term은 주어진 reward model에 샘플들이 overfit 되는 것을 방지하며, beta는 그 세기를 결정하는 hyperparameter이다. 
5. 위 목적 함수를 최대화하는 target distribution은 다음과 같이 정리된다. 
% 
6. 하지만 target distribution에서 직접 샘플링을 하는 것은 분모의 normalization constant로 인해 불가능하다. 
7. 다음 섹션에서는 기존 diffusion model을 활용한 이전 연구들이 particle sampler을 활용해서 target distribution로부터 샘플을 생성한 방식에 대해 살펴보겠다. 
% 
% 
}


% We consider a task that aims to generate a sample from a pretrained flow model that maximizes the reward. 
% We denote 

% Inference-time scaling~\cite{} aims to improve the performance of pretrained models by allocating additional computation during inference. 
% In this work, we focus on aligning a flow model to a given reward model, such as one that measures text alignment~\cite{} or aesthetic quality~\cite{}. 

% Formally, our objective is to generate samples $\B{x}_0 \in \mathbb{R}^d$ from a pretrained flow model while maximizing a given reward model $r(\cdot): \mathbb{R}^{d} \rightarrow \mathbb{R}$, without deviating significantly from $p_\theta$, the original distribution of the pretrained model. Following previous works on reward alignment~\cite{On Reinforcement Learning and Distribution Matching for Fine-Tuning Language Models with No Catastrophic Forgetting, Fine-tuning Continuous, Bridging Model-Based Optimization and Generative Modeling via Conservative Fine-Tuning of Diffusion Models}, this objective can be formulated as follows:
% \begin{align}
%     \label{eq:reward_max_obj}
%     \nonumber
%     p^{*}(\B{x}) &= \argmax_{p} \  \mathbb{E}_{\B{x} \sim p} \underbrace{\left[ r(\B{x}) \right]}_{\text{Reward}} -\beta \underbrace{\mathcal{D}_{\text{KL}} \left[ p(\B{x}) | p_{\text{ref}}(\B{x}) \right]}_{\text{KL Regularization}} \\
%     \label{eq:target_distribution}
%     &= \frac{1}{Z} p_\theta(\B{x}) \exp \left( \frac{r(\B{x})}{\beta} \right),
% \end{align}
% where $Z$ is a normalization constant, and $\beta$ is a regularization hyperparameter (proof in Appendix~\ref{sec:appendix_optimal_marginal_target}). 

% To sample from the target distribution in Eq.~\ref{eq:target_distribution}, previous works utilizing diffusion models employed particle sampling, which will be discussed in the next section. 

\vspace{-0.4\baselineskip}
\subsection{Particle Sampling Using Diffusion Model}
\vspace{-0.5\baselineskip}
\label{sec:motivation}

A pretrained diffusion model generates data by drawing an initial sample from the standard Gaussian distribution and iteratively sampling from the learned conditional distribution $p_{\theta}(\B{x}_{t-\Delta t} | \B{x}_{t})$. 
Building on this, previous works~\cite{Levine:2018RL, Uehara:2024Bridging} have shown that data from the target distribution in Eq.~\ref{eq:target_distribution} can be generated by performing the same denoising process while replacing the conditional distribution $p_{\theta}(\B{x}_{t-\Delta t} | \B{x}_{t})$ with the \emph{optimal policy}:

{\small
    \begin{align}
    \label{eq:optimal_policy}
        p^{*}_\theta(\B{x}_{t-\Delta t} | \B{x}_{t}) = \frac{ p_\theta(\B{x}_{t-\Delta t} | \B{x}_{t}) \exp \left(\frac{v(\B{x}_{t-\Delta t})}{\beta}  \right)}{\int p_\theta(\B{x}_{t-\Delta t} | \B{x}_{t}) \exp \Bigl(\frac{v(\B{x}_{t-\Delta t})}{\beta} \Bigr) \mathrm{d}\B{x}_{t-\Delta t}},
    \end{align}
}%
where the details are presented in Appendix~\ref{sec:appendix_optimal_soft_policy}. 
We denote $v(\cdot): \mathbb{R}^d \rightarrow \mathbb{R}$ as the optimal value function that estimates the expected future reward of the generated samples at current timestep. 
Following previous works~\cite{Chung:2023DPS,Kim:2025DAS, Li2024:SVDD,Bansal:2023UGD}, we approximate the value function using the posterior mean computed via Tweedie’s formula~\cite{Robbins1992}, given by $v(\B{x}_t) \approx r(\B{x}_{0|t})$, where $\B{x}_{0|t} \coloneq \mathbb{E}_{\B{x}_0 \sim p_{\theta}(\B{x}_0 | \B{x}_t)} \left[ \B{x}_0 \right]$. 

\mycomment{
1. Diffusion model은 조건부 확률 $p_{\text{ref}}(\B{x}_{t-\Delta t} | \B{x}_{t})$로부터 샘플을 반복적으로 추출하여 데이터를 생성한다. 
2. 이러한 점에서 착안하여 이전 연구들은 optimal policy에서 샘플을 반복적으로 추출함으로써 target distribution의 데이터를 생성해내는 것을 보였다. 

}

% Previous works~\cite{} have shown that one can generate a sample from the target distribution in Eq.~\ref{eq:target_distribution} by iteratively sampling from the \textit{optimal policy}, which is formulated as:
% \begin{align}
% \label{eq:optimal_policy}
%     p^{*}(\B{x}_{t-\Delta t} | \B{x}_{t}) &= \frac{p_{\text{ref}}(\B{x}_{t-\Delta t} | \B{x}_{t}) \exp \left(\frac{\hat{r}(\B{x}_{t-\Delta t})}{\beta}  \right)}{\int p_{\text{ref}}(\B{x}_{t-\Delta t} | \B{x}_{t}) \exp \Bigl(\frac{\hat{r}(\B{x}_{t-\Delta t})}{\beta} \Bigr) \mathrb{d}\B{x}_{t-\Delta t}},
% \end{align}
% where $p_{\text{ref}}(\B{x}_{t-\Delta t} | \B{x}_{t})$ is the reverse process of a pretrained diffusion model, and $\hat{r}(\cdot): \mathbb{R}^d \rightarrow \mathbb{R}$ is a soft value function that estimates the expected future reward of the generated samples. 
% Following previous works~\cite{DPS, DAS, SVDD, Universal guidance}, the value function is approximated through the posterior mean $\B{x}_{0|t}$. 
% We present detailed derivations in Appendix.~\ref{sec:appendix_optimal_soft_policy}. 



% Previous works on diffusion models seek to sample $\B{x}_0$ from the target distribution by iteratively drawing samples from the \textit{optimal policy}~\cite{}:
% \begin{align}
% \label{eq:optimal_policy}
%     p^{*}(\B{x}_{t-\Delta t} | \B{x}_{t}) &= \frac{p_{\text{ref}}(\B{x}_{t-\Delta t} | \B{x}_{t}) \exp \left(\frac{\hat{r}(\B{x}_{t-\Delta t})}{\beta}  \right)}{\int p_{\text{ref}}(\B{x}_{t-\Delta t} | \B{x}_{t}) \exp \Bigl(\frac{\hat{r}(\B{x}_{t-\Delta t})}{\beta} \Bigr) \mathrb{d}\B{x}_{t-\Delta t}},
% \end{align}
% where $\hat{r}(\cdot): \mathbb{R}^d \rightarrow \mathbb{R}$ is a soft value function that estimates the expected future reward of the generated samples. 
% Following previous works~\cite{DPS, DAS, SVDD, Universal guidance}, the value function is approximated through the posterior mean $\B{x}_{0|t}$. We present detailed proof in Appendix~\ref{sec:appendix_optimal_soft_policy}. 

Since directly sampling from the optimal policy distribution in Eq.~\ref{eq:optimal_policy} is nontrivial, one can first approximate the distribution using importance sampling while taking $p_{\theta}(\B{x}_{t-\Delta t} | \B{x}_{t})$ as the \emph{proposal distribution}:
%previous works have relied on \textit{particle sampling}~\cite{}. 
% In the context of particle sampling, we consider $p_{\theta}(\B{x}_{t-\Delta t} | \B{x}_{t})$ as the proposal distribution, which is stochastic by design in diffusion models, allowing multiple different particles $\B{x}_{t-\Delta t}$ to be sampled conditioned on $\B{x}_t$. 
% Each step of particle sampling for diffusion models involves importance sampling: 
{\small
\begin{align}
    \label{eq:importance_sampling}
    p^{*}_\theta(\B{x}_{t-\Delta t} | \B{x}_{t}) \approx \sum_{i=1}^K \frac{w^{(i)}_{t-\Delta t}}{\sum_{j=1}^K w^{(j)}_{t-\Delta t}} \delta_{\B{x}^{(i)}_{t-\Delta t}}, \quad
    % \nonumber
    \{ \B{x}^{(i)}_{t-\Delta t} \}_{i=1}^{K} \sim p_\theta(\B{x}_{t-\Delta t} | \B{x}_{t}), 
    % & \quad w^{(i)}_{t-\Delta t} = \exp( \frac{r ( \B{x}^{(i)}_{0|t-\Delta t} )}{\beta} ),
\end{align}
}%
where $K$ is the number of particles, $w^{(i)}_{t-\Delta t} = \exp( v ( \B{x}^{(i)}_{t-\Delta t} ) / \beta )$ is the weight, and $\delta_{\B{x}^{(i)}_{t-\Delta t}}$ is a Dirac distribution. 
SVDD~\cite{Li2024:SVDD} proposed an approximate sampling method for the optimal policy by selecting the sample with the largest weight from Eq.~\ref{eq:importance_sampling}. 
% On the other hand, \citet{Kim:2025DAS} employ a different probabilistic approach, drawing particles from a multinomial distribution. 
% denotes the empirical distribution over the sampled particles. 
%At each denoising step, the sample with the largest importance weight is selected from the fixed set of particles and propagated through the denoising process~\cite{}.

% Notably, a key factor in sampling high-reward $\B{x}_{t-\Delta t}$ using importance sampling is designing the proposal distribution to sufficiently cover the distribution of high-reward samples. 
Notably, a key factor in seeking high-reward samples using particle sampling is defining the proposal distribution to sufficiently cover the distribution of high-reward samples.
% Consider a scenario where high-reward samples reside in a low density region of the original data distribution. 
Consider a scenario where high-reward samples reside in a low density region of the original data distribution, which is common when generating complex or highly specific samples that deviate from the mode of the pretrained model distribution. 
In this case, the proposal distribution must have a sufficiently large variance to effectively explore these low density regions.
This highlights the importance of the \textit{stochasticity} of the proposal distribution, which has been instrumental in the successful adoption of particle sampling in diffusion models. 
In contrast, flow models~\cite{Lipman:2023CFM} employ a \textit{deterministic} sampling process, where all particles $\B{x}_{t-\Delta t}$ drawn from $\B{x}_{t}$ are identical. 
This restricts the applicability of particle sampling methods in flow models. 
One of the main contributions is the investigation of how these particle sampling methods can be efficiently applied to flow models.

%Inspired by the success of particle sampling in diffusion models, we propose inference-time approaches to expand the exploration space and facilitate the use of particle sampling methods in flow models. 
To this end, we propose an inference-time approach that introduces stochasticity into the generative process of flow models to enable particle sampling. 
We first transform the deterministic sampling process of flow models into a stochastic process (Sec.~\ref{sec:ode_sde_conversion}). 
We further identify a sampling trajectory that expands the search space of the flow models (Sec.~\ref{sec:scheduler_conversion}). 
Note that while stochastic sampling and trajectory conversion have been studied in prior works, their primary goals have been to improve sample quality~\cite{xu2023restart, yeo2025stochsync, salimans2022progressive, kingma2021variational, Ma2024:sit} or to accelerate inference~\cite{holderrieth2024generator, shaul2023bespoke, shaul2023kinetic, Karras2022:EDM}. To the best of our knowledge, we are the first to investigate sampling stochasticity and trajectory conversion for efficient particle-based sampling in flow models.

% Additionally, previous particle sampling methods in diffusion models employed fixed compute (number of particles) across all denoising timesteps, potentially limiting exploration. 
Additionally, previous particle sampling methods in diffusion models allocated a fixed computational budget (i.e., a uniform number of particles) across all denoising timesteps, potentially limiting exploration.
We explore sampling with the rollover strategy, which adaptively allocates compute across timesteps during the sampling process (Sec.~\ref{sec:roll_over}). 

% Additionally, we introduce a scheduler conversion that converts the sample trajectory of flow models into that of diffusion models, further expanding the exploration space (Sec.\ref{sec:scheduler_conversion}). 
% Lastly, we explore adaptive scheduling of particle size along the temporal axis to efficiently assign the computing resources (Sec.~\ref{sec:temporal_axis_scaling}).

% To overcome this limitation, we propose an inference-time approach that transforms the deterministic sampling process of flow models into a stochastic process (Sec.\ref{sec:ode_sde_conversion}). 
% Additionally, we introduce a scheduler conversion that converts the sample trajectory of flow models into that of diffusion models, further expanding the exploration space (Sec.\ref{sec:scheduler_conversion}). 
% Lastly, we adopt adaptive particle-size sampling, which dynamically allocates computational resources across denoising steps to enhance efficiency (Sec.~\ref{sec:temporal_axis_scaling}).


% \Jaihoon{TODO} However, applying particle sampling to flow models, which utilize a \textit{deterministic} sampling approach, presents inherent limitations. 
% \begin{remark}[Degeneracy in Dirac delta proposals]
%     \label{remark:deterministic_proposal}
%     When the proposal distribution $p_\theta(\B{x}_{t-\Delta t} | \B{x}_t)$ is a Dirac delta distribution, the effective sample size of the drawn particles collapses to 1, thereby restricting exploration of low-density regions where potential optimal samples may reside. 
% \end{remark}

% In Sec.~\ref{sec:method}, we propose a novel inference-time, approach that enables the application of efficient search algorithms~\cite{} originally developed for diffusion models to flow models. 

% To overcome this, we introduce post-training SDE conversion, which converts ODE sampling process into an SDE, allowing the application of search algorithms that have demonstrated strong performance in diffusion model inference-time scaling to flow models. 
% While SDE conversion expands the search space, we further propose post-training scheduler conversion to enhance exploration. Lastly, inspired by existing search algorithms that scale with a fixed number of particles at each denoising step, we introduce our method, which adapts the scaling behavior along the temporal axis.

\mycomment{
1. LLM과 Diffusion literature에서는 inference-time의 computation을 늘림으로써 사용자의 preference에 보다 더 align된 샘플을 생성할 수 있는 inference-time reward alignment가 많은 주목을 받았고, 다양한 방법론들이 소개되었다. 
2. 하지만 최근 빠른 생성 시간과 high-fidelity로 각광을 받고있는 flow 모델에 대해서는 이러한 연구가 상대적으로 큰 관심을 받지 못하였다. 
3. 본 연구에서는 기존 LLM과 Diffusion 모델에서 큰 주목을 받고 있는 inference-time scaling을 flow 모델에 대해서 적용시키고자 한다. 
% 
4. 이전 연구에서 제시하였듯이(On Reinforcement Learning and Distribution Matching for Fine-Tuning Language Models with no Catastrophic For- getting, Fine-tuning continuous), 우리는 i) 주어진 reward를 최대화시키면서 ii) 기존의 분포에서 크게 벗어나지 않는 sample을 생성하는 것이 목표이다. 이러한 objective는 다음과 같이 표기된다: 
5. 다음 section에서는 diffusion 모델의 literature에서 위 분포를 만족시키는 샘플을 inference-time에 추출하는 particle sampling에 대하여 소개하겠다. 
% 
% 
% 
6. 기존 Diffusion 모델을 활용한 연구에서는 위 분포에서 샘플링을 하기 위해 여러 개의 particle을 기반으로 한 Monte Carlo 기반의 방식들이 소개되었다. 
7. (MC 샘플링 예시). 이러한 방식을 활용한 Diffusion 모델들이 성공적인 alignemnt를 수행할 수 있었던 이유에는 proposal 분포 p(x_{t-\Delta t}|x_{t})의 stochasticity가 더 넓은 exploration space를 제공해주었기 때문이다. 
8. 하지만 deterministic한 샘플링 방식을 사용하는 flow 모델에 대해서는 이러한 particle 샘플링 방식을 적용하는데 한계가 존재한다. (증명)
9. 따라서 우리는 Diffusion 모델의 inference-time scaling에 좋은 성능을 보인 search algorithm들을 flow 모델에서도 마찬가지로 적용시키기 위해 기존 ODE를 SDE로 변환하는 post-training SDE conversion을 소개한다. 비록 SDE conversion이 더 넓은 search space 가져오는데 도움을 주지만 더 넓은 exploration을 하기 위해 우리는 post-training scheduler conversion을 제안한다. 마지막으로 우리는 기존의 search algorithm들은 매 denoising step에서 같은 수의 particle에 대해서만 scaling을 하는데 inspire하여 temporal axis의 scaling을 adaptive하게 바꾼 our method를 소개한다. 
}





% -----------------------------------------------------------

% \section{Problem Definition}
% \label{sec:objective}
% In this section, we define our objective of generating samples $\B{x}_0 \in \mathbb{R}^d$ that maximize a given reward model $r(\cdot): \mathbb{R}^{d} \rightarrow \mathbb{R}$. Our goal is to produce high-reward samples without deviating too far from the original distribution of the pretrained model $p_\theta$. Following~\cite{Fine-Tuning of Continuous-Time Diffusion Models as Entropy-Regularized Control, DPO}, this problem can be formulated as an entropy-regularized reward maximization objective:
% \begin{align}
%     \label{eq:reward_max_obj}
%     p^{*}(\B{x}) &= \argmax_{p} \  \mathbb{E}_{\B{x} \sim p} \underbrace{\left[ r(\B{x}) \right]}_{\text{Reward}} -\beta \underbrace{\mathcal{D}_{\text{KL}} \left[ p(\B{x}) | p_{\text{ref}}(\B{x}) \right]}_{\text{Regularization}} \\
%     \label{eq:target_distribution}
%     &= \frac{1}{Z} p_\theta(\B{x}) \exp \left( \frac{r(\B{x})}{\beta} \right),
% \end{align}
% where $Z$ is a normalization constant, and $\beta$ is a regularization hyperparameter (proof in Appendix~\ref{sec:appendix_optimal_marginal_target}). 

% To sample $\B{x}_0$ from the target distribution $p^{*}$, we leverage the optimal policy $p^{*}(\B{x}_{t-\Delta t} | \B{x}_{t})$~\cite{Fine-Tuning of Continuous-Time Diffusion Models as Entropy-Regularized Control, Bridging Model-Based Optimization and Generative Modeling via Conservative Fine-Tuning of Diffusion Models
% }:
% \begin{align}
% \label{eq:optimal_policy}
%     p^{*}(\B{x}_{t-\Delta t} | \B{x}_{t}) &= \frac{p_{\text{ref}}(\B{x}_{t-\Delta t} | \B{x}_{t}) \exp \left(\frac{\hat{r}(\B{x}_{t-\Delta t})}{\beta}  \right)}{\int p_{\text{ref}}(\B{x}_{t-\Delta t} | \B{x}_{t}) \exp \Bigl(\frac{\hat{r}(\B{x}_{t-\Delta t})}{\beta} \Bigr) \mathrb{d}\B{x}_{t-\Delta t}},
% \end{align}
% where $\hat{r}(\cdot): \mathbb{R}^d \rightarrow \mathbb{R}$ is a soft value function that estimates the expected future reward at $t=0$. Following previous works~\cite{DPS, DAS, SVDD, Universal guidance}, the value function is approximated through the posterior mean $\B{x}_{0|t}$. A detailed proof is provided in Appendix~\ref{sec:appendix_optimal_soft_policy}.

% However, sampling from the optimal policy in Eq.~\ref{eq:optimal_policy} is not tractable due to the normalization constant. 
% In the next section, we introduce training-free, search approaches designed to address this challenge. 
